\documentclass{article}

\usepackage{amsmath}
\usepackage{amssymb}
\usepackage{geometry}
\geometry{margin=2.5cm}
\title{F42: Fortran coding guides}
\author{Asis Hallab, ``el Bobito''}

\begin{document}

\maketitle

\section{Memory Passing Behavior in Fortran (F42 Standard for Tensor Omics)}

\subsection{General Principles}

In Fortran (standards 77--2018), arguments to subroutines and functions are, by default, passed \textbf{by reference} rather than by copying.  
This applies to scalars, arrays, and derived types.

Thus:

\begin{itemize}
    \item Scalars (e.g., \texttt{integer}, \texttt{real}) are passed by reference.
    \item Arrays (e.g., \texttt{integer, dimension(:)}, \texttt{real, dimension(:,:)}) are passed by reference.
    \item Derived types are passed by reference.
\end{itemize}

\subsection{When Copies Are Made}

Despite the default pass-by-reference behavior, hidden copies may be generated by the compiler in the following situations:

\begin{itemize}
    \item \textbf{Passing expressions}: e.g., \texttt{call process(a+b)} requires evaluation and temporary storage.
    \item \textbf{Passing strided array sections}: e.g., \texttt{array(1:10:2)} is non-contiguous and may trigger copying.
    \item \textbf{Passing non-contiguous slices}: particularly relevant for rows in a 2D array.
\end{itemize}

\subsection{Tensor Omics DataTables: Column and Row Passing}

Tensor Omics DataTables store typed data in 2D Fortran arrays of the form:

\begin{itemize}
    \item \texttt{real\_array(n\_rows, n\_real\_columns)}
    \item \texttt{int\_array(n\_rows, n\_int\_columns)}
    \item \texttt{char\_array(n\_rows, n\_char\_columns)}
\end{itemize}

\textbf{Accessing Columns}:

\begin{itemize}
    \item Example: \texttt{real\_array(:, j)}
    \item Behavior: Passed by reference, \textbf{no copy made}.
    \item Reason: Slicing across the first index in column-major layout yields contiguous memory.
\end{itemize}

\textbf{Accessing Rows}:

\begin{itemize}
    \item Example: \texttt{real\_array(i, :)}
    \item Behavior: Typically triggers a hidden copy.
    \item Reason: Slicing across the second index accesses non-contiguous memory addresses, requiring packing into a contiguous temporary.
\end{itemize}

\subsection{Summary Tables}

\subsubsection{General Fortran Behavior}

\begin{center}
\begin{tabular}{|l|l|l|}
\hline
\textbf{Action} & \textbf{Example} & \textbf{Copy Made?} \\
\hline
Pass full array & \texttt{array} & No \\
Pass contiguous slice & \texttt{array(:, j)} & No \\
Pass strided slice & \texttt{array(1:10:2, j)} & Yes \\
Pass expression & \texttt{array(:, j) * 2.0} & Yes \\
\hline
\end{tabular}
\end{center}

\subsubsection{Tensor Omics DataTables Specifics}

\begin{center}
\begin{tabular}{|l|l|l|}
\hline
\textbf{Access Type} & \textbf{Example} & \textbf{Copy Made?} \\
\hline
Access full column & \texttt{real\_array(:, j)} & No \\
Access full row & \texttt{real\_array(i, :)} & Yes \\
\hline
\end{tabular}
\end{center}

\subsection{Best Practice Recommendation for Tensor Omics}

\begin{itemize}
    \item Always design core functions to work with columns if possible.
    \item Avoid passing strided slices or expressions directly to subroutines.
    \item Treat rows carefully: copy manually if absolutely needed, otherwise prefer per-element processing.
\end{itemize}

\subsection{Conclusion}

Efficient memory handling is critical for Tensor Omics.  
Understanding Fortran's passing conventions allows maximization of speed, minimization of hidden temporaries, and preservation of F42 computational standards.

\end{document}
